\documentclass[aps,prc,onecolumn,superscriptaddress,nofootinbib,10pt]{revtex4-2}
\usepackage{amsmath,amssymb}
\usepackage{booktabs}
\usepackage[colorlinks=true,linkcolor=blue,citecolor=blue,urlcolor=blue]{hyperref}

\newcommand{\nB}{n_B}
\newcommand{\nC}{n_C}
\newcommand{\nS}{n_S}
\newcommand{\muB}{\mu_B}
\newcommand{\muC}{\mu_C}
\newcommand{\muS}{\mu_S}
\newcommand{\mutB}{\tilde\mu_B}
\newcommand{\phH}{\mathrm{H}}
\newcommand{\phQ}{\mathrm{Q}}

\begin{document}

\title{The hadron--quark mixed phase: a one-parameter family between the
       Gibbs and Maxwell constructions \\
       --- as implemented in \texttt{eos/mixed}}

\author{M.~Guerrini}
\affiliation{University of Ferrara}

\begin{abstract}
The \texttt{eos/mixed} subpackage builds the first-order deconfinement
transition between a hadronic phase and a quark phase. It is a
\emph{composite engine}, not a model: it supplies no matter of its own, and
couples two validated bulk equations of state --- DD2~\cite{Typel2010,
HempelSchaffnerBielich2010} for hadrons and vMIT~\cite{Chodos1974,Gomes2019}
for quarks --- through a single narrow interface, the phase-adapter contract.
A continuous parameter $\eta \in [0,1]$ selects how much of electric-charge
neutrality is imposed phase by phase rather than on the mixture as a whole,
so that $\eta = 0$ is the Gibbs construction~\cite{Glendenning1992},
$\eta = 1$ is Maxwell, and intermediate values interpolate between them in
the sense of Constantinou et al.~\cite{Constantinou2023}. This note states
the contract, the equilibrium conditions, how the unknown vector and the
residual are \emph{derived} from a per-charge regime declaration rather than
enumerated per mode, and what a mixed table reports beyond a point equation
of state. It is written from the code, as its specification: every equation
here is asserted by the \texttt{verify/} suite or by \texttt{test/mixed/}.
\end{abstract}

\maketitle

\section{What is being solved}

At a given total baryon density $\nB$, temperature $T$ and construction
parameter $\eta$, a hadronic phase H occupying volume fraction $1-\chi$
coexists with a quark phase Q occupying $\chi$. The two phases are separate
thermodynamic systems, each with its own composition and its own set of
conserved-charge chemical potentials; coexistence is the statement that
certain of those potentials are shared and that certain charges are conserved
only in the volume average. Which ones is exactly the content of the
construction, and it is the only thing that changes between $\eta = 0$ and
$\eta = 1$.

$\chi$ is an unknown of the solve, not an input, and it is deliberately
\emph{not} clamped to $[0,1]$. Its value classifies the density: $\chi \le 0$
means the point is still pure hadronic, $\chi \ge 1$ that it is already pure
quark, and $0 < \chi < 1$ that it lies inside the coexistence window. The
onset and offset densities are located by exactly that sign change, which is
why the solver must be allowed to return the analytic continuation outside
the window.

\section{The phase-adapter contract}

Everything this engine knows about DD2 and about vMIT passes through one
interface. An adapter is a map
%
\begin{equation}
\big( \text{the phase's conserved-charge potentials},\, T \big)
\;\longmapsto\;
\big\{ n_i \big\},\; \nB,\, \nC,\, \nS,\; P,\, \varepsilon,\, s;\;
\muB,\, \muC,\, \muS ,
\label{eq:adapter}
\end{equation}
%
which solves that phase's own internal self-consistency --- the meson field
equations for a density-dependent relativistic mean field, the vector-field
fixed point for a bag model --- at the \emph{given} potentials, and reports
the resulting block. There is no $\eta$ in Eq.~(\ref{eq:adapter}), no mixing
and no neutrality: those are conditions on the \emph{pair} of phases. An
adapter describes one phase in isolation, and pairing a different hadronic or
quark model is writing an adapter, not editing the solver.

Two properties of Eq.~(\ref{eq:adapter}) are load-bearing.

\paragraph{One decomposition.} Both adapters report the same projection
%
\begin{equation}
\mu_i = B_i \muB + C_i \muC + S_i \muS ,
\label{eq:projection}
\end{equation}
%
with $C$ the electric charge of strongly-interacting matter only (leptons
excluded) and $S = +1$ per $s$ quark, the opposite of the PDG sign. Because
both phases speak Eq.~(\ref{eq:projection}), the coexistence conditions can
be written in terms of $(\muB, \muC, \muS)$ without either side knowing which
engine produced the other.

\paragraph{Determinism.} An adapter must be a function of its arguments
alone: the same potentials must give the same block, to the last digit,
however the point was reached. The mixed residual is differentiated
numerically, so an adapter that remembered the previous trial point would
make the Jacobian differentiate that memory along with the physics. This is
why warm starts are passed in as explicit arguments and why the one thing
that \emph{is} cached --- the hadronic phase's starting field configuration
--- is a constant of the solve rather than a running state.

\paragraph{The hadronic potential is the kinetic one.} The hadronic adapter
takes $\mutB = \muB - \Sigma^R$, with $\Sigma^R$ the rearrangement
self-energy of the density-dependent couplings. Since $\Sigma^R$ is itself a
function of the phase density, which is an unknown of the phase-internal
solve, carrying $\muB$ would put that circularity inside the outer iteration;
carrying $\mutB$ leaves it inside the adapter, where it belongs. The physical
$\muB$ is restored at assembly and is what the matching condition uses.

\section{The $\eta$ family of constructions}

Write the electron (and, when enabled, muon) gas as two populations. A
\emph{local} population lives inside the structures and neutralises its own
phase; a \emph{global} population is spread over the mixture and neutralises
only the average. They carry weights $\eta$ and $1-\eta$:
%
\begin{equation}
\text{local: } \mu_e^{\mathrm{L},\phH},\ \mu_e^{\mathrm{L},\phQ}
\quad \text{(weight } \eta),
\qquad
\text{global: } \mu_e^{\mathrm{G}}
\quad \text{(weight } 1-\eta).
\label{eq:leptonsplit}
\end{equation}
%
Neutrality is then imposed twice, each where its population lives:
%
\begin{align}
\text{local } (\eta > 0): &\qquad
   \nC^\phH = n_\ell^\phH , \qquad \nC^\phQ = n_\ell^\phQ ,
\label{eq:neutral_local} \\
\text{global } (\eta < 1): &\qquad
   (1-\chi)\, \nC^\phH + \chi\, \nC^\phQ = n_\ell^{\mathrm{G}} ,
\label{eq:neutral_global}
\end{align}
%
where $n_\ell$ is the total negatively-charged lepton density of that domain.
At $\eta = 1$ only Eq.~(\ref{eq:neutral_local}) survives: each phase is
separately neutral, no charge is exchanged between them, and the two phases
can only coexist at one pressure --- the Maxwell plateau, with a genuine
density jump. At $\eta = 0$ only Eq.~(\ref{eq:neutral_global}) survives: each
phase may be charged, charge is exchanged freely, and the pressure rises
continuously through the window --- Gibbs. In between, both populations
exist. Physically the intermediate case stands in for the finite surface
tension and Coulomb energy of the mixed-phase structures, which suppress
charge separation without forbidding it; here it is a controlled
interpolation, not a derivation from a surface tension.

Mechanical equilibrium carries the same weighting. Only the local leptons sit
inside the structures whose pressures must balance, so
%
\begin{equation}
P^\phH + \eta\, P_\ell^\phH \;=\; P^\phQ + \eta\, P_\ell^\phQ .
\label{eq:mechanical}
\end{equation}
%
The global leptons, the photons and any trapped neutrinos are common to both
phases and cancel from Eq.~(\ref{eq:mechanical}) identically.

\section{Conserved charges as a declaration}

Each conserved quantity in $\{B, C, S, L_e\}$ is treated in one of three ways:

\begin{description}
\item[\textsc{global}] the potential is shared, $\mu_j^\phH = \mu_j^\phQ$,
  and the charge is conserved in the volume average,
  $(1-\chi) n_j^\phH + \chi\, n_j^\phQ = Y_j \nB$;
\item[\textsc{local}] the potential is per-phase and the charge is conserved
  inside each phase separately;
\item[\textsc{not conserved}] the charge is not conserved at all: its
  potential is eliminated from the unknown vector.
\end{description}

$B$ is \textsc{global} in every mode --- $\muB^\phH = \muB^\phQ$ is what makes
the phases coexist at all. The four named equilibrium modes are four
assignments of the remaining three charges:

\begin{center}
\begin{tabular}{lccccl}
\toprule
mode & $B$ & $C$ & $S$ & $L_e$ & independent variables \\
\midrule
\texttt{beta\_eq\_neutrinoless}     & global & --     & --     & --     & $(\nB, T)$ \\
\texttt{beta\_eq\_neutrino\_trapped}& global & --     & --     & global & $(\nB, Y_{L_e}, T)$ \\
\texttt{fixed\_YC}                  & global & global & --     & --     & $(\nB, Y_C, T)$ \\
\texttt{fixed\_YC\_YS}              & global & global & global & --     & $(\nB, Y_C, Y_S, T)$ \\
\bottomrule
\end{tabular}
\end{center}

\noindent
(``--'' is \textsc{not conserved}.) Nothing else in the engine enumerates
modes. The unknown vector, the residual list and the analytic Jacobian are all
\emph{assembled by reading the regimes}, so the four named modes are four
configurations of one solver and an unnamed combination needs no new code.

Two combinations are refused rather than mis-assembled: per-phase strangeness
($S$ \textsc{local}) is not wired, and $L_e$ \textsc{local} is not a defined
mode at all --- the neutrino mean free path is far larger than the
mixed-phase structures, so neutrinos cannot be localised in one of them.

\section{The equilibrium system}

Because the adapters have already absorbed each phase's internal
self-consistency, the unknowns here are only conserved-charge potentials,
$\chi$, and the split lepton potentials:
%
\begin{equation}
x = \Big(
\underbrace{\mutB^\phH,\ \muB^\phQ,\ \chi}_{\text{always}} ,\;
\underbrace{\muC^\phH, \muC^\phQ \ \text{or}\ \muC}_{C \text{ global}} ,\;
\underbrace{\muS}_{S \text{ global}} ,\;
\underbrace{\mu_{\nu_e}}_{L_e \text{ global}} ,\;
\underbrace{\mu_e^{\mathrm{L},\phH}, \mu_e^{\mathrm{L},\phQ}}_{\eta > 0} ,\;
\underbrace{\mu_e^{\mathrm{G}}}_{\eta < 1}
\Big) ,
\label{eq:unknowns}
\end{equation}
%
four to nine numbers, rather than the full per-species $(\mu_i, n_i)$ vector
a direct formulation would carry. The residual rows are, in order:
%
\begin{align}
\muB^\phH - \muB^\phQ &= 0 ,
&& \text{(baryon potentials match)}
\label{eq:res_mub} \\
(1-\chi)\, \nB^\phH + \chi\, \nB^\phQ - \nB &= 0 ,
&& \text{(the average reproduces the target density)}
\label{eq:res_nb} \\
P^\phH + \eta P_\ell^\phH - P^\phQ - \eta P_\ell^\phQ &= 0 ,
&& \text{(mechanical equilibrium, Eq.~\ref{eq:mechanical})}
\label{eq:res_P}
\end{align}
%
followed, by regime, by
%
\begin{align}
(1-\chi)\, \nC^\phH + \chi\, \nC^\phQ - Y_C \nB &= 0 ,
&& C\ \textsc{global}
\label{eq:res_C} \\
\muC^\phH + \eta\, \mu_e^{\mathrm{L},\phH}
   - \muC^\phQ - \eta\, \mu_e^{\mathrm{L},\phQ} &= 0 ,
&& C\ \textsc{global with leptons}
\label{eq:res_Cmatch} \\
(1-\chi)\, \nS^\phH + \chi\, \nS^\phQ - Y_S \nB &= 0 ,
&& S\ \textsc{global}
\label{eq:res_S} \\
\eta \big[ (1-\chi) n_e^\phH + \chi n_e^\phQ \big]
   + (1-\eta) n_e^{\mathrm{G}} + n_{\nu_e} - Y_{L_e} \nB &= 0 ,
&& L_e\ \textsc{global}
\label{eq:res_L}
\end{align}
%
and closed by the neutrality rows Eqs.~(\ref{eq:neutral_local}) and
(\ref{eq:neutral_global}), each present exactly when its lepton population is.

The charge-matching row Eq.~(\ref{eq:res_Cmatch}) is where the construction
enters the charge sector: at $\eta = 0$ it reduces to $\muC^\phH = \muC^\phQ$,
the Gibbs statement that the two phases share a charge potential, and at
$\eta = 1$ the local electron potentials shift it by exactly the amount that
separate neutrality demands. The global electron potential does not appear in
it, because it neutralises the average rather than either phase.

In $\beta$ equilibrium ($C$ \textsc{not conserved}) the charge potential is
not an unknown at all: it is eliminated by the weak condition applied with the
$\eta$-weighted electron potential of that phase,
%
\begin{equation}
\muC^{I} = \mu_{\nu_e}
   - \big[ \eta\, \mu_e^{\mathrm{L},I} + (1-\eta)\, \mu_e^{\mathrm{G}} \big] ,
\qquad I \in \{\phH, \phQ\} ,
\label{eq:beta}
\end{equation}
%
which for transparent matter ($\mu_{\nu_e} = 0$) is the repository's sign
convention $\muC + \mu_e = 0$, i.e.\ $\muC = \mu_p - \mu_n$. Strangeness
self-equilibrates in that case, $\muS = 0$.

\section{Thermodynamics of the mixture}

The matter phases are volume-averaged; the local leptons carry weight $\eta$
and are themselves volume-averaged; the global leptons carry weight $1-\eta$;
photons and trapped neutrinos are uniform across the whole mixture and are
counted once:
%
\begin{align}
P &= P^\phH + \eta P_\ell^\phH + (1-\eta) P^{\mathrm{G}}_\ell
     + P_{\nu} + P_{\gamma} ,
\label{eq:Ptot} \\
\varepsilon &= (1-\chi)\, \varepsilon^\phH + \chi\, \varepsilon^\phQ
     + \eta \big\langle \varepsilon_\ell \big\rangle
     + (1-\eta)\, \varepsilon^{\mathrm{G}}_\ell
     + \varepsilon_{\nu} + \varepsilon_{\gamma} ,
\label{eq:epstot} \\
s &= (1-\chi)\, s^\phH + \chi\, s^\phQ
     + \eta \big\langle s_\ell \big\rangle + (1-\eta)\, s^{\mathrm{G}}_\ell
     + s_{\nu} + s_{\gamma} ,
\label{eq:stot}
\end{align}
%
with $\langle X_\ell \rangle = (1-\chi) X_\ell^\phH + \chi X_\ell^\phQ$. The
pressure is read off one phase rather than averaged because
Eq.~(\ref{eq:res_P}) has made the two equal.

\paragraph{Euler / Hugenholtz--Van Hove.} The identity
%
\begin{equation}
\varepsilon + P = T s + \sum_i \mu_i n_i
\label{eq:euler}
\end{equation}
%
must hold for the mixture as a whole, with the sum weighted exactly as
Eqs.~(\ref{eq:epstot})--(\ref{eq:stot}) weight $\varepsilon$ and $s$, and with
$\mu_\gamma = 0$. Unlike the pure-phase case this is not an algebraic
identity: it holds only if every thermal and lepton term has been weighted
consistently, so it is the sharpest single check on the assembly and it is
asserted at $10^{-8}$ relative on every solved point.

\section{The two sound speeds}

A first-order transition has two, and which is physical depends on how fast
the matter is disturbed relative to the rate at which one phase converts into
the other.

The \emph{equilibrium} speed $c_{\mathrm{eq}}^2 = dP/d\varepsilon$ is taken
along the solved sequence, with $\chi$ free to readjust: a compression is
answered by converting hadrons into quarks rather than by raising the
pressure. Through a Maxwell window it therefore vanishes identically, and
through a Gibbs window it dips but stays finite. This is the speed that enters
the TOV equations and whose causality bound $0 \le c^2 \le 1$ is checked
before any table is integrated.

The \emph{frozen} (adiabatic) speed $c_{\mathrm{ad}}^2$ holds $\chi$ fixed.
The mixture is compressed faster than the phases can convert, the pressure has
to rise, and $c_{\mathrm{ad}}$ does not collapse in the window. Freezing
$\chi$ is the part that matters: freezing only the charge fractions would let
the solve readjust $\chi$ and return to the plateau. In addition each phase is
compressed by the same factor, each keeps its own $Y_C$ and $Y_S$, and the
leptons are re-neutralised against the frozen total charge. That convention is
stated in full in the module documentation, because a different one gives a
different number; with $\chi = 0$ and the leptons switched off it reduces
exactly to the pure-hadronic adiabatic sound speed of \texttt{eos/dd2}.

The gap between the two is what drives a composition $g$-mode across the
transition.

\section{What a mixed table reports}

A composite engine returns more than a point equation of state. Alongside the
rows --- bulk thermodynamics, $\chi$, the phase label, and every conserved
charge resolved by phase --- a table returns the \emph{windows}: the onset and
offset densities located on each isotherm and each combination of the fixed
fractions. Those are what the transition curves $\nB^{\mathrm{onset}}(T)$ and
$\nB^{\mathrm{offset}}(T)$ are made of, and they are part of the result rather
than something the caller recovers afterwards by scanning the rows for where
$\chi$ crossed $0$ and $1$.

The per-phase decomposition satisfies three cheap invariants at every solved
point,
%
\begin{equation}
Y_B^\phH + Y_B^\phQ = 1 , \qquad
Y_C^\phH + Y_C^\phQ = Y_C , \qquad
Y_S^\phH + Y_S^\phQ = Y_S ,
\label{eq:partition}
\end{equation}
%
on the volume-weighted convention $Y_j^I = w_I n_j^I / \nB$ with
$w_\phH = 1-\chi$ and $w_\phQ = \chi$. The partition is what the global sums
cannot show: at fixed total $Y_C$ the hadronic phase can be far more
positively charged than the average while the quark phase carries the
compensating negative charge, and how far that separation goes is exactly what
$\eta$ controls.

\section{Numerics}

\paragraph{Scaling and the gate.} Every residual row is made dimensionless
before it is judged: density rows are divided by $\nB$, potential equalities
by a fixed potential scale, and the mechanical row by the larger of the two
phase pressures. A state is accepted when the largest scaled residual is below
$10^{-10}$. The rows carry mixed units --- densities of order
$10^{-1}\,\mathrm{fm}^{-3}$ against potentials of order
$10^{3}\,\mathrm{MeV}$ --- so a gate on the raw norm would be dominated by
whichever row happens to be largest and would accept states satisfying the
others only loosely.

\paragraph{Warm starts.} These dominate the cost, because the adapters are
called once per residual evaluation. Three levels are used, in order of what
they buy: the hadronic phase's starting field configuration is a constant of
the solve and is computed once; each density seeds the next along the stiff
$\nB$ axis; and each isotherm seeds the next. On top of that the mixed system
is solved only \emph{inside} the located window --- outside it the far cheaper
pure-phase solvers answer the same question --- and the window search on one
isotherm is told where the previous isotherm found its boundaries, since they
move smoothly with $T$. A hint that produces a boundary on the edge of the
hinted range is discarded and the full search repeated, so a window that
genuinely disappears is still reported as gone.

\paragraph{Jacobian.} The reference path uses the solver's own numeric
Jacobian and is the correctness oracle. A hand-assembled analytic Jacobian is
available as the fast path, validated against finite differences and required
to reach the same root; where a trial point drives a phase solve out of its
domain, both paths answer with a large penalty residual so the outer iteration
backs off instead of aborting.

\paragraph{Non-convergence.} Reported, never raised, at the public boundary:
the engine is used inside parameter scans that walk into regions where no
hybrid equation of state exists, and a scan must be able to score such a point
and move on. A converged point outside $[0,1]$ in $\chi$ is not a failure ---
it is how the engine says which side of the transition the density is on.

\begin{acknowledgments}
Written as the specification of \texttt{eos/mixed}.
\end{acknowledgments}

\bibliography{../../docs/eos}

\end{document}
